%% mwart Polish article starter — mwart class (part of the MW bundle).
%%
%% The mwart class itself is NOT included in this download: it ships with
%% every modern TeX distribution (and with latextemplates.com's compiler),
%% so bundling a copy would only risk pinning you to an older version than
%% your own system already has.
%%
%% Class documentation: https://ctan.org/pkg/mwcls
%% Licence: LPPL-1.2 (or later), (c) 2000-2025 Marcin Woli\'nski, verified
%% directly in the shipped mwart.cls header.
%%
%% mwart is a Polish-oriented alternative to the standard article class
%% (part of the MW bundle, alongside mwbk/mwrep), with typesetting
%% conventions tuned for Polish typography. It is language-agnostic on its
%% own -- \usepackage[polish]{babel} + \usepackage{polski} below are what
%% actually localise section names, hyphenation and quotation marks.
\documentclass[a4paper,11pt]{mwart}

\usepackage[T1]{fontenc}
\usepackage[utf8]{inputenc}
\usepackage[polish]{babel}
\usepackage{polski}

\title{Tytuł artykułu}
\author{Ada Lovelace}
\date{\today}

\begin{document}

\maketitle

\begin{abstract}
Zastąp ten tekst streszczeniem. Dobre streszczenie opisuje problem, co
zostało zrobione i jakie są wyniki --- w tej kolejności, bez odwołań do
bibliografii.
\end{abstract}

\section{Wprowadzenie}
Zacznij tutaj. Cytuj wcześniejsze prace w ten sposób~\cite{knuth1984texbook}.

\section{Metoda}
\label{sec:metoda}
Opisz, co zostało zrobione.

\section{Wyniki}
\label{sec:wyniki}
Przedstaw wyniki.

\section{Podsumowanie}
Podsumuj, co to oznacza.

\bibliographystyle{plain}
\bibliography{references}

\end{document}
